% Abstract and Zusammenfassung. Written 2026-09-28 on Abram's instruction, replacing the
%   template's lorem ipsum, and trimmed twice the same day so that each block sits on one page.
%   Snapshots, oldest first: abstract.preWritten.bak.tex (the template's lorem ipsum),
%   abstract.preTrim.bak.tex (the first write, 453 w English and 467 w German, two pages each),
%   abstract.preDeFix.bak.tex (English at its final 387 w, German still at 381 w and still
%   overflowing).
%
% ⚠⚠ THE GERMAN BLOCK IS WRAPPED IN otherlanguage{ngerman}, AND THIS IS NOT COSMETIC.
%   softlang.cls loads babel as \RequirePackage[ngerman,english]{babel} (line 85) when the class
%   option `german` is NOT set, which is this thesis's case, so english is the main language and
%   ngerman is merely available. Without the switch the Zusammenfassung is typeset with ENGLISH
%   hyphenation patterns, and a German compound -- Formelrepräsentation, Relevanzbewerter,
%   Formelübereinstimmung, ARQMath-3-Kollektion -- then cannot be broken at all. Every one of them
%   is forced whole onto a line, the justification goes ragged, and the block runs long for
%   reasons that have nothing to do with how much it says. That was the overflow Abram reported at
%   09:45, after the word count had already been brought under the English one.
%   The class also loads hyphsubst with ngerman-x-latest (line 90) precisely for compounds; that
%   too only takes effect once the language is ngerman.
%   ⚠ ngerman makes " an active character (see the build log). The block contains no " and must
%   not acquire one -- use \enquote or German quotes if quoting ever becomes necessary here.
%
% TWO PAGES, his instruction: \clearpage between the two blocks. The document is `oneside`, so
%   \clearpage inserts no blank page and ausarbeitung.tex's own \cleardoublepage after the \input
%   still carries the reader into the table of contents. The German block comes first, as the
%   template had it; swapping the two is one cut and paste if the English should lead.
% The template's instruction line -- "Gemäß der Prüfungsordnung ist ein Abstract in deutscher und
%   englischer Sprache verpflichtend." -- is removed: it addresses the author, not the reader.
%
% ⚠ THE PAGE BUDGET IS MEASURED, NOT GUESSED. The fullest solid-prose pages of the build hold
%   443, 433 and 422 words (pages 26, 29, 31); a \subsection* heading costs roughly three and a
%   half lines, about 40 words, and \parskip is zero in `book`. That puts the ceiling near 400
%   words for English. German sets fewer words per line -- 5.9 characters per word against 5.0 --
%   so the German block is held well under that ceiling rather than to the same number.
%   THE HISTORY, so that no one restores a longer German block believing the ceiling is 400:
%   453 w (two pages) -> 381 w (still two) -> 356 w with the ngerman fix below (still one sentence
%   over, reported 09:49) -> the present 330. English has been at its final 387 since the first
%   trim and has always fitted. The three German cuts came out of the opening paragraph each time,
%   which is the only part of either block carrying no result.
%
% NO NUMBER FROM THE REGISTRY. The only figures are the two containment counts, 70 and 30 of the
%   78 topics (manifest §14 R7, OPT, fanout, H+M), which need no scale to be read, and the 78
%   topics themselves. Every nDCG' value in this thesis is computed under the judge's labels and
%   means nothing beside a number from anywhere else (§4.4), so none appears here; an abstract is
%   the one place a bare 0.551 would be read as a benchmark figure. Nothing here can therefore go
%   stale against the registry.
%
% BANS HONOURED: no hypothesis ID, no RQ number, no "significant", no ratio or magnitude claim, no
%   first person, no section reference, no citation, "formulas" not "formulae". No coined term --
%   *blocked*, *flat*, *formula block* and *floor* are described rather than named, since their
%   definition sites (§5.2.2, §6.2) sit thirty to fifty pages further on. The merge sentence claims
%   only that the arrangement decides the outcome, never that it matters more than the
%   representation choice (§1c-i decision 1).
%
% SCOPE. The abstract names mathematical questions and the collection, and claims nothing about
%   scientific papers -- §1.1's concession, and it sidesteps the undecided title (§15 O-6).
%
% ⚠ THREE THINGS NO FURTHER TRIM MAY UNDO.
%   (i) "Adding all three" CARRIES ITS QUALIFIER, *when each votes separately* / *wenn jeder für
%       sich abstimmt*. Without it the claim is false -- it holds only of the arrangement where
%       each formula ranking votes on its own (§5.4.4), and the blocked system adds all three as
%       well (§5.4.5), so the unqualified form contradicts the merge sentence two lines below it.
%   (ii) The formula-search sentence stays: it is the answer to "was the component simply
%       broken?" (§5.3), and it is the one sentence an examiner will look for.
%   (iii) Both containment counts stay. 70 without 30 is the opposite finding.
%   Where the next cut goes, if a block still spills: the opening paragraph, the only part of
%   either block carrying no result.
%
% ⚠ DASHES: --- in the English block, -- in the German, which is the German convention. Not an
%   inconsistency.
% ⚠ THE GERMAN IS A TRANSLATION AND ITS TERMS ARE CHOICES, not settled vocabulary: sparse and
%   dense read as "dünn besetzt" and "dicht"; the ARQMath assessors as "Annotatoren", to keep
%   "Gutachter" for the examiners of the title page; "answer retrieval" as "Antwortsuche", beside
%   "Formelsuche". Sweep all four together if any one of them changes.

\begin{otherlanguage}{ngerman}
\subsection*{Zusammenfassung}

Eine mathematische Frage wird in Worten und Formeln zugleich gestellt, und ein Retrievalsystem
muss beides lesen. In einem gewöhnlichen Textindex steht eine Formel nur als die Zeichenfolge,
mit der sie geschrieben wurde; die Anordnung, in der ihre Bedeutung liegt, fehlt dort. Die
Struktur zu indexieren hilft bei der Formelsuche; ob es auch bei der Antwortsuche hilft, ist eine
andere Frage, denn relevant ist eine Antwort, wenn sie die Frage löst, nicht, wenn ihre Formeln
übereinstimmen.

Diese Arbeit baut drei strukturbewusste Formelindizes, einen je Formelrepräsentation, neben einem
dünn besetzten und einem dichten Textindex und führt ihre Ranglisten in zehn Anordnungen
zusammen. Die fünfzehn Konfigurationen werden auf der Antwortsuche über die 78 Topics der
ARQMath-3-Kollektion bewertet, unter einem zuvor an den menschlichen Urteilen validierten
Relevanzbewerter, und jede liefert anschließend einen Kontext an einen Generator unter einem
festen Prompt.

Formelevidenz verbessert die Antwortsuche nicht. Kommt einer der drei Formelindizes zum stärksten
reinen Textretrieval hinzu, bleibt das System dahinter zurück, ebenso bei allen dreien mit je
eigener Stimme. Wie sie zusammengeführt werden, entscheidet über den Abstand: Bündelt man die
drei Formelranglisten zu einer Stimme, zieht das System mit dem reinen Textretrieval gleich, und
nichts hier übertrifft es. Diese Anordnung verdient ihren Platz, indem sie die Formelevidenz
daran hindert zu schaden, nicht indem sie liefert, was den Textranglisten fehlt. Am Index liegt
es nicht; bei der Formelsuche schneidet er vergleichbar mit einem veröffentlichten System ab. Die
Formeln eines Topics stecken auf 70 der 78 Topics stärker in den relevanten als in den nicht
relevanten Antworten, unter den am besten passenden aber, wo sich eine Rangliste entscheidet, nur
noch auf 30. Formelübereinstimmung findet die richtige Nachbarschaft und kann innerhalb dieser
nicht ordnen.

Ein abgerufener Kontext erreicht die erzeugte Antwort dennoch, die faktisch korrekter ausfällt
als eine ohne Kontext geschriebene. Wie viel einer Antwort aus ihrem Kontext stammt und wie viel
davon richtig ist, steigt mit der Qualität des Retrievals dahinter, wenn auch nicht gleichmäßig
und nicht bei jeder Größe des Generators. Das Retrieval erreicht die Antwort; es schreibt aber
nur wenig davon.
\end{otherlanguage}

\clearpage

\subsection*{Abstract}

A mathematical question is asked in words and formulas together, and a retriever has to read
both. A formula enters an ordinary text index as the characters it was typed with. Rename the
variables and the index sees an unrelated string. What the index never holds is the arrangement
itself, and that is where a formula's meaning lies. Indexing that structure is what formula
search does; whether it also helps when answers are what is retrieved is a separate question,
because an answer is relevant when it resolves the question and not when its formulas match.

This thesis builds three structure-aware formula indexes, one per formula representation, beside
a sparse and a dense text index, and merges the rankings they return in ten arrangements. Fifteen
configurations in all are scored on answer retrieval over the 78 topics of the ARQMath-3
collection, under a relevance judge validated beforehand against the human assessors, and each of
them then supplies a context to one generator under one fixed prompt.

Formula evidence does not improve answer retrieval. Adding any one of the three indexes to the
strongest text-only retrieval leaves the system behind that retrieval alone, and so does adding
all three when each of them votes separately. How the rankings are merged decides how far behind:
collapsing the three formula rankings into a single vote before fusing draws level with text-only
retrieval, and nothing built here passes it. That arrangement earns its place by keeping formula
evidence from doing harm rather than by supplying what the text rankings lack. The failure is not
the index, which performs comparably with a published system on formula search. A topic's
formulas are contained in the answers judged relevant more than in those judged not relevant on
70 of the 78 topics, but among the best-matching answers, where a ranking is decided, on only 30.
Formula matching finds the right neighbourhood and cannot order inside it.

A retrieved context nonetheless reaches the generated answer, which is more factually correct
than one written from no context at all. Both how much of an answer comes from its context and
how much of it is correct rise with the quality of the retrieval behind it, though not equally
and not at every size of generator. Retrieval reaches the answer; it does not write much of it.
